\documentclass[10pt,a4paper]{article}
\usepackage[T1]{fontenc}
\usepackage[utf8]{inputenc}
\usepackage[margin=2cm,landscape]{geometry}
\usepackage{amsmath,amssymb}
\usepackage{array,booktabs,longtable}
\usepackage{xcolor}
\usepackage{enumitem}
\usepackage[hidelinks]{hyperref}
\newcommand{\lean}[1]{\textcolor{blue!50!black}{\texttt{\detokenize{#1}}}}
\setlength{\LTpre}{6pt}\setlength{\LTpost}{6pt}
\renewcommand{\arraystretch}{1.25}
\begin{document}


\begin{center}{\Large\bfseries Proposition 5: the lower bound $4/3$ for the edit distance}\\[4pt]

{\small Correspondence between the paper and the Lean~4/Mathlib formalisation}\end{center}

\noindent\textbf{Paper.} J.~Breitung, T.~Gervens, M.~Grohe. \emph{Graph Similarity Meets Optimal Transport} (file \texttt{Graph\_Similarity\_Meets\_Optimal\_Transport.pdf} in the repository). Statement, equation, page and line numbers refer to this PDF; \texttt{docs/NUMBERING.md} lists the Lean files and declarations of all numbered statements.\par\smallskip

The folder \texttt{GsimOT/Proposition5\_EditLowerBound} formalises Proposition~5 (p.~6) with the construction of Appendix~C.1, Lemmas~23--25 and Claim~26, and Theorem~3~(1). Shared notions and lemmas (the edit distance, Observation~21, Lemma~22) come from \texttt{GsimOT/Common}; they are listed at the end.\par\paragraph{Checking.} Run \texttt{bash build.sh} in the repository root (see \texttt{README.md}). The script builds the library and runs \texttt{scripts/CheckAxioms.lean}, which must report only the standard axioms \texttt{propext}, \texttt{Classical.choice}, \texttt{Quot.sound}: no \texttt{sorry} and no additional axioms.

\paragraph{Main statement in Lean.} \lean{GsimOT.Prop5.proposition5_dEdit}: for every odd $N\ge3$, the graphs $G=$ \lean{GN N} and
$H=$ \lean{HN N} on the common vertex set \lean{VN N} satisfy
\[ d_{\mathrm{ed}}(G,H)\ge4N-2,\qquad d^{\mathrm{OT}}_{\mathrm{ed}}(G,H)\le3N+1; \]
\lean{GsimOT.Prop5.proposition5} states it for the costs ($C_{\mathrm{ed}}(\pi)\ge4N-2$ for all
bijections $\pi$, $C_{\mathrm{ed}}(Q)\le3N+1$ for some $Q\in\mathrm{DS}(V,W)$).
\lean{GsimOT.Prop5.proposition5_delta} (and \lean{GsimOT.Prop5.theorem3_item1}, Theorem~3~(1)): for every
$\varepsilon>0$ there is an odd $N\ge3$ with
$\delta_{\mathrm{ed}}(G,H)\ge(\frac43-\varepsilon)\,\delta^{\mathrm{OT}}_{\mathrm{ed}}(G,H)$; by
\lean{GsimOT.Prop5.proposition5_delta_of_le} this holds for every odd $N\ge3$ with
$\frac{10}{3(3N+1)}\le\varepsilon$.

\section*{Correspondence: Proposition 5 and its proof}

\begin{longtable}{>{\raggedright\arraybackslash}p{3.2cm} >{\centering\arraybackslash}p{1.1cm} >{\raggedright\arraybackslash}p{8.4cm} >{\raggedright\arraybackslash}p{10.2cm}}
\toprule
\textbf{Paper} & \textbf{Page} & \textbf{Statement} & \textbf{Lean declaration(s) and location} \\
\midrule\endhead
\bottomrule\endfoot
{}Appendix~C.1 (construction, l.~677--687) & 21 & Markers $z_0,\dots,z_{N-1}$; cores $S_G=\{v_0,v_2,\dots,v_{2N-2}\}$, $S_H=\{w_1,\dots,w_{2N-1}\}$ at their positions in $\mathbb Z_{2N}$; distinct core vertices adjacent iff their positions sum to $0$ mod $2N$; $z_r$ adjacent to the core vertices on $I_r=\{(r+s)\bmod 2N\mid s<N\}$ & \lean{GsimOT.Prop5.Base}\newline{\scriptsize\texttt{Proposition5\_EditLowerBound/Construction.lean:37}} \newline \lean{GsimOT.Prop5.semicircle}\newline{\scriptsize\texttt{Proposition5\_EditLowerBound/Construction.lean:40}} \newline \lean{GsimOT.Prop5.mem_semicircle}\newline{\scriptsize\texttt{Proposition5\_EditLowerBound/Construction.lean:43}} \newline \lean{GsimOT.Prop5.posG}\newline{\scriptsize\texttt{Proposition5\_EditLowerBound/Construction.lean:53}} \newline \lean{GsimOT.Prop5.posH}\newline{\scriptsize\texttt{Proposition5\_EditLowerBound/Construction.lean:56}} \newline \lean{GsimOT.Prop5.coreGraph}\newline{\scriptsize\texttt{Proposition5\_EditLowerBound/Construction.lean:68}} \newline \lean{GsimOT.Prop5.Gp}\newline{\scriptsize\texttt{Proposition5\_EditLowerBound/Construction.lean:88}} \newline \lean{GsimOT.Prop5.Hp}\newline{\scriptsize\texttt{Proposition5\_EditLowerBound/Construction.lean:92}} \\ \midrule
{}Appendix~C.1 (construction, l.~692--694) & 22 & For odd $N$: core edges $\{c,c'\}$ with $c+c'=N$ in $G'$ and with $c+c'+1=N$, $c\ne c'$ in $H'$; $v_0$ and $w_N$ are isolated & \lean{GsimOT.Prop5.Gp_adj_inl_inl}\newline{\scriptsize\texttt{Proposition5\_EditLowerBound/Construction.lean:106}} \newline \lean{GsimOT.Prop5.Hp_adj_inl_inl}\newline{\scriptsize\texttt{Proposition5\_EditLowerBound/Construction.lean:117}} \\ \midrule
{}Appendix~C.1 (construction, l.~688--691) & 21--22 & $G,H$ by Lemma~22, markers ordered by their degree in $G'$, $\Delta=N+1$, $\tau=4N-1$, $\delta=1$; all degrees are at most $N+1$, and the degrees of $z_r$ in $G'$ and $H'$ differ by one & \lean{GsimOT.Prop5.markerEnum}\newline{\scriptsize\texttt{Proposition5\_EditLowerBound/Construction.lean:249}} \newline \lean{GsimOT.Prop5.markerEnum_mono}\newline{\scriptsize\texttt{Proposition5\_EditLowerBound/Construction.lean:254}} \newline \lean{GsimOT.Prop5.deg_Gp_le}\newline{\scriptsize\texttt{Proposition5\_EditLowerBound/Construction.lean:218}} \newline \lean{GsimOT.Prop5.deg_Hp_le}\newline{\scriptsize\texttt{Proposition5\_EditLowerBound/Construction.lean:231}} \newline \lean{GsimOT.Prop5.deg_Gp_inr}\newline{\scriptsize\texttt{Proposition5\_EditLowerBound/Construction.lean:189}} \newline \lean{GsimOT.Prop5.deg_Hp_inr}\newline{\scriptsize\texttt{Proposition5\_EditLowerBound/Construction.lean:200}} \newline \lean{GsimOT.Prop5.markerEnum_deg}\newline{\scriptsize\texttt{Proposition5\_EditLowerBound/Construction.lean:262}} \newline \lean{GsimOT.Prop5.mK}\newline{\scriptsize\texttt{Proposition5\_EditLowerBound/Construction.lean:270}} \newline \lean{GsimOT.Prop5.GN}\newline{\scriptsize\texttt{Proposition5\_EditLowerBound/Construction.lean:276}} \newline \lean{GsimOT.Prop5.HN}\newline{\scriptsize\texttt{Proposition5\_EditLowerBound/Construction.lean:279}} \\ \midrule
{}Lemma~23, Fig.~4 (l.~696--699, 703--706) & 22 & $\ell(v_i,w_j)=\min_k|i-j+2kN|$; mapping $v_i$ to $w_j$ changes the adjacency to exactly $\ell(v_i,w_j)$ markers & \lean{GsimOT.Prop5.cycDist}\newline{\scriptsize\texttt{Proposition5\_EditLowerBound/Lemma23\_MarkerCost.lean:37}} \newline \lean{GsimOT.Prop5.ell}\newline{\scriptsize\texttt{Proposition5\_EditLowerBound/Lemma23\_MarkerCost.lean:73}} \newline \lean{GsimOT.Prop5.card_semicircle_ne}\newline{\scriptsize\texttt{Proposition5\_EditLowerBound/Lemma23\_MarkerCost.lean:51}} \newline \lean{GsimOT.Prop5.sum_marker_inl}\newline{\scriptsize\texttt{Proposition5\_EditLowerBound/Lemma23\_MarkerCost.lean:77}} \newline \lean{GsimOT.Prop5.sum_marker_inr}\newline{\scriptsize\texttt{Proposition5\_EditLowerBound/Lemma23\_MarkerCost.lean:82}} \\ \midrule
{}Lemma~23 (l.~700--702) & 22 & For compliant $Q$, the core-marker pairs contribute exactly $2\sum Q(v_i,w_j)\,\ell(v_i,w_j)$ to $C_{\mathrm{ed}}(Q)$ & \lean{GsimOT.Prop5.extMarkers}\newline{\scriptsize\texttt{Proposition5\_EditLowerBound/Lemma23\_MarkerCost.lean:90}} \newline \lean{GsimOT.Prop5.coreCost}\newline{\scriptsize\texttt{Proposition5\_EditLowerBound/Lemma23\_MarkerCost.lean:120}} \newline \lean{GsimOT.Prop5.costEdit_extMarkers}\newline{\scriptsize\texttt{Proposition5\_EditLowerBound/Lemma23\_MarkerCost.lean:127}} \newline \lean{GsimOT.Prop5.lemma23}\newline{\scriptsize\texttt{Proposition5\_EditLowerBound/Lemma23\_MarkerCost.lean:165}} \\ \midrule
{}Proof of Lemma~24 (l.~708--714) & 22 & $\pi_\pm(v_i)=w_{i\pm1}$, $Q=\frac12(P_{\pi_-}+P_{\pi_+})$; $\ell(v_i,\pi_\pm(v_i))=1$ & \lean{GsimOT.Prop5.predCore}\newline{\scriptsize\texttt{Proposition5\_EditLowerBound/Lemma24\_FractionalCost.lean:37}} \newline \lean{GsimOT.Prop5.avgPerm}\newline{\scriptsize\texttt{Proposition5\_EditLowerBound/CoreEdges.lean:170}} \newline \lean{GsimOT.Prop5.Rwit}\newline{\scriptsize\texttt{Proposition5\_EditLowerBound/Lemma24\_FractionalCost.lean:106}} \newline \lean{GsimOT.Prop5.Qwit}\newline{\scriptsize\texttt{Proposition5\_EditLowerBound/Lemma24\_FractionalCost.lean:109}} \newline \lean{GsimOT.Prop5.ell_self}\newline{\scriptsize\texttt{Proposition5\_EditLowerBound/Lemma24\_FractionalCost.lean:116}} \newline \lean{GsimOT.Prop5.ell_predCore}\newline{\scriptsize\texttt{Proposition5\_EditLowerBound/Lemma24\_FractionalCost.lean:122}} \\ \midrule
{}Proof of Lemma~24, Fig.~5 (l.~715--724) & 22--23 & Core-core pairs: the core edges are preserved iff one endpoint moves forward and the other backward, except near $v_0$ and $w_N$; the Lean computes their cost as $\frac14(2(N-1)+2+2+2(N-1))=N$ & \lean{GsimOT.Prop5.pres_eq}\newline{\scriptsize\texttt{Proposition5\_EditLowerBound/CoreEdges.lean:135}} \newline \lean{GsimOT.Prop5.pres_id_id}\newline{\scriptsize\texttt{Proposition5\_EditLowerBound/Lemma24\_FractionalCost.lean:57}} \newline \lean{GsimOT.Prop5.pres_pred_pred}\newline{\scriptsize\texttt{Proposition5\_EditLowerBound/Lemma24\_FractionalCost.lean:65}} \newline \lean{GsimOT.Prop5.pres_id_pred}\newline{\scriptsize\texttt{Proposition5\_EditLowerBound/Lemma24\_FractionalCost.lean:75}} \newline \lean{GsimOT.Prop5.pres_pred_id}\newline{\scriptsize\texttt{Proposition5\_EditLowerBound/Lemma24\_FractionalCost.lean:90}} \newline \lean{GsimOT.Prop5.coreCost_avgPerm}\newline{\scriptsize\texttt{Proposition5\_EditLowerBound/CoreEdges.lean:210}} \newline \lean{GsimOT.Prop5.coreCost_witness}\newline{\scriptsize\texttt{Proposition5\_EditLowerBound/Lemma24\_FractionalCost.lean:128}} \\ \midrule
{}Lemma~24 (l.~707, 725) & 22--23 & $d^{\mathrm{OT}}_{\mathrm{ed}}(G,H)\le3N+1$, from $C_{\mathrm{ed}}(Q)=3N$ for the witness $Q$ & \lean{GsimOT.Prop5.costEdit_witness}\newline{\scriptsize\texttt{Proposition5\_EditLowerBound/Lemma24\_FractionalCost.lean:139}} \newline \lean{GsimOT.Prop5.lemma24}\newline{\scriptsize\texttt{Proposition5\_EditLowerBound/Lemma24\_FractionalCost.lean:145}} \\ \midrule
{}Proof of Lemma~25 (l.~727--729) & 23 & If $C_{\mathrm{ed}}(\pi)<\tau$, then $\pi$ fixes the markers and maps $S_G$ onto $S_H$ (Lemma~22); the pairs involving the clique only add to the cost & \lean{GsimOT.lemma22_bij_edit}\newline{\scriptsize\texttt{Common/Lemma22\_LockingMarkersEdit.lean:41}} \newline \lean{GsimOT.Prop5.restrictInl}\newline{\scriptsize\texttt{Proposition5\_EditLowerBound/Lemma25\_BijectionCost.lean:44}} \newline \lean{GsimOT.Prop5.costEditBij_restrict}\newline{\scriptsize\texttt{Proposition5\_EditLowerBound/Lemma25\_BijectionCost.lean:56}} \newline \lean{GsimOT.Prop5.costEditBij_sumCongr}\newline{\scriptsize\texttt{Proposition5\_EditLowerBound/Lemma25\_BijectionCost.lean:76}} \\ \midrule
{}Proof of Lemma~25, Claim~26 (l.~729--742) & 23--24 & Displacements $d_i\in\{-N+2,\dots,N\}$ with $\ell(v_i,\pi(v_i))=|d_i|$; $\sum d_i\equiv N\pmod{2N}$, so $|\sum d_i|\ge N$ & \lean{GsimOT.Prop5.disp}\newline{\scriptsize\texttt{Proposition5\_EditLowerBound/Lemma25\_BijectionCost.lean:86}} \newline \lean{GsimOT.Prop5.ell_eq_abs_disp}\newline{\scriptsize\texttt{Proposition5\_EditLowerBound/Lemma25\_BijectionCost.lean:92}} \newline \lean{GsimOT.Prop5.one_le_ell}\newline{\scriptsize\texttt{Proposition5\_EditLowerBound/Lemma25\_BijectionCost.lean:99}} \newline \lean{GsimOT.Prop5.sum_disp}\newline{\scriptsize\texttt{Proposition5\_EditLowerBound/Lemma25\_BijectionCost.lean:115}} \newline \lean{GsimOT.Prop5.sum_disp_ge}\newline{\scriptsize\texttt{Proposition5\_EditLowerBound/Lemma25\_BijectionCost.lean:125}} \\ \midrule
{}Claim~26 (l.~737--749) & 24 & $U\ge N+2k$; preserved edges have $d_i+d_{i'}=0$ & \lean{GsimOT.Prop5.Preserved}\newline{\scriptsize\texttt{Proposition5\_EditLowerBound/Lemma25\_BijectionCost.lean:136}} \newline \lean{GsimOT.Prop5.disp_add_disp_partner}\newline{\scriptsize\texttt{Proposition5\_EditLowerBound/Lemma25\_BijectionCost.lean:154}} \newline \lean{GsimOT.Prop5.claim26}\newline{\scriptsize\texttt{Proposition5\_EditLowerBound/Lemma25\_BijectionCost.lean:167}} \\ \midrule
{}Proof of Lemma~25 (l.~750--751) & 24 & Both cores have $(N-1)/2$ edges; their ordered pairs contribute $2(N-1-2k)$ & \lean{GsimOT.Prop5.sum_adj_Gp_core}\newline{\scriptsize\texttt{Proposition5\_EditLowerBound/CoreEdges.lean:96}} \newline \lean{GsimOT.Prop5.sum_adj_Hp_core}\newline{\scriptsize\texttt{Proposition5\_EditLowerBound/CoreEdges.lean:108}} \newline \lean{GsimOT.Prop5.mixCore}\newline{\scriptsize\texttt{Proposition5\_EditLowerBound/CoreEdges.lean:125}} \newline \lean{GsimOT.Prop5.pres}\newline{\scriptsize\texttt{Proposition5\_EditLowerBound/CoreEdges.lean:130}} \newline \lean{GsimOT.Prop5.mixCore_eq}\newline{\scriptsize\texttt{Proposition5\_EditLowerBound/CoreEdges.lean:145}} \\ \midrule
{}Lemma~25 (l.~726, 753) & 23--24 & $C_{\mathrm{ed}}(\pi)\ge4N-2$ for every $\pi$, i.e.\ $d_{\mathrm{ed}}(G,H)\ge4N-2$ & \lean{GsimOT.Prop5.lemma25_bij}\newline{\scriptsize\texttt{Proposition5\_EditLowerBound/Lemma25\_BijectionCost.lean:237}} \newline \lean{GsimOT.Prop5.lemma25}\newline{\scriptsize\texttt{Proposition5\_EditLowerBound/Lemma25\_BijectionCost.lean:275}} \\ \midrule
{}Proposition~5 (l.~221--222) & 6 & $d_{\mathrm{ed}}(G,H)\ge4N-2$ and $d^{\mathrm{OT}}_{\mathrm{ed}}(G,H)\le3N+1$; in terms of costs: $C_{\mathrm{ed}}(\pi)\ge4N-2$ for all $\pi$, $C_{\mathrm{ed}}(Q)\le3N+1$ for some $Q$ & \lean{GsimOT.Prop5.proposition5_dEdit}\newline{\scriptsize\texttt{Proposition5\_EditLowerBound/Proposition5.lean:31}} \newline \lean{GsimOT.Prop5.proposition5}\newline{\scriptsize\texttt{Proposition5\_EditLowerBound/Proposition5.lean:38}} \\ \midrule
{}Proof of Proposition~5 (l.~754--758) & 24 & $d_{\mathrm{ed}}\ge4N-2=(\frac43-\frac{10}{3(3N+1)})(3N+1)\ge(\frac43-\frac{10}{3(3N+1)})\,d^{\mathrm{OT}}_{\mathrm{ed}}$; same normalisation; choose $N$ with $\frac{10}{3(3N+1)}\le\varepsilon$ & \lean{GsimOT.Prop5.proposition5_delta_of_le}\newline{\scriptsize\texttt{Proposition5\_EditLowerBound/Proposition5.lean:45}} \\ \midrule
{}Proposition~5 (l.~222--223) & 6 & For every $\varepsilon>0$ there is an $N$ with $\delta_{\mathrm{ed}}(G,H)\ge(\frac43-\varepsilon)\,\delta^{\mathrm{OT}}_{\mathrm{ed}}(G,H)$ & \lean{GsimOT.Prop5.proposition5_delta}\newline{\scriptsize\texttt{Proposition5\_EditLowerBound/Proposition5.lean:68}} \\ \midrule
{}Theorem~3 (1) (l.~124) & 4 & For every $\varepsilon>0$ some $G,H$ have $\delta_{\mathrm{ed}}\ge(\frac43-\varepsilon)\,\delta^{\mathrm{OT}}_{\mathrm{ed}}$ & \lean{GsimOT.Prop5.theorem3_item1}\newline{\scriptsize\texttt{Proposition5\_EditLowerBound/Proposition5.lean:82}} \\
\end{longtable}

\section*{Shared notions and lemmas used (from \texttt{GsimOT/Common})}

\begin{longtable}{>{\raggedright\arraybackslash}p{3.2cm} >{\centering\arraybackslash}p{1.1cm} >{\raggedright\arraybackslash}p{8.4cm} >{\raggedright\arraybackslash}p{10.2cm}}
\toprule
\textbf{Paper} & \textbf{Page} & \textbf{Statement} & \textbf{Lean declaration(s) and location} \\
\midrule\endhead
\bottomrule\endfoot
{}Section~1, Section~2 & 2, 5 & $C_{\mathrm{ed}}(Q)=\sum_{v,w,v',w'}|D(v,w,v',w')|\,Q(v,w)\,Q(v',w')$ and $C_{\mathrm{ed}}(\pi)=C_{\mathrm{ed}}(P_\pi)=\|A_G^\pi-A_H\|_{\mathrm{ed}}$ & \lean{GsimOT.costEdit}\newline{\scriptsize\texttt{Common/EditDistance.lean:35}} \newline \lean{GsimOT.costEditBij}\newline{\scriptsize\texttt{Common/EditDistance.lean:39}} \newline \lean{GsimOT.costEditBij_eq}\newline{\scriptsize\texttt{Common/EditDistance.lean:60}} \\ \midrule
{}Section~1, Eq.~(1), Section~2, Appendix~A & 2, 5, 14--15 & $d_{\mathrm{ed}}=\min_\pi C_{\mathrm{ed}}(\pi)$ and $d^{\mathrm{OT}}_{\mathrm{ed}}=\min_{Q\in\mathrm{DS}(V,W)}C_{\mathrm{ed}}(Q)$ (infimum, attained); $d^{\mathrm{OT}}_{\mathrm{ed}}\le d_{\mathrm{ed}}$ & \lean{GsimOT.dEdit}\newline{\scriptsize\texttt{Common/EditDistance.lean:42}} \newline \lean{GsimOT.dOTEdit}\newline{\scriptsize\texttt{Common/EditDistance.lean:45}} \newline \lean{GsimOT.exists_dOTEdit_eq}\newline{\scriptsize\texttt{Common/EditDistance.lean:109}} \newline \lean{GsimOT.dOTEdit_le_dEdit}\newline{\scriptsize\texttt{Common/EditDistance.lean:102}} \\ \midrule
{}Section~1, Section~2 & 1--2, 5 & $\delta_{\mathrm{ed}}=n^{-2}d_{\mathrm{ed}}$, $\delta^{\mathrm{OT}}_{\mathrm{ed}}=n^{-2}d^{\mathrm{OT}}_{\mathrm{ed}}$ (scaling by $1/n^2$) & \lean{GsimOT.deltaEdit}\newline{\scriptsize\texttt{Common/EditDistance.lean:48}} \newline \lean{GsimOT.deltaOTEdit}\newline{\scriptsize\texttt{Common/EditDistance.lean:51}} \\ \midrule
{}Section~1, Section~2, Appendix~A & 2, 5, 14 & $C_{\mathrm{ed}}(Q)=\|B(G,H,Q)\|_{\mathrm{ed}}$ for $Q\ge0$ (entrywise $1$-norm $\|\cdot\|_{\mathrm{ed}}=\|\cdot\|_{(1)}$) & \lean{GsimOT.entrywiseNorm}\newline{\scriptsize\texttt{Common/MatrixNorms.lean:53}} \newline \lean{GsimOT.costEdit_eq_entrywiseNorm_B}\newline{\scriptsize\texttt{Common/MatrixNorms.lean:119}} \\ \midrule
{}Section~2, Eq.~(3) & 5 & Birkhoff--von Neumann: $Q=\sum_i\lambda_iP_{\pi_i}$ with $\lambda_i>0$, $\sum_i\lambda_i=1$ (also: over all $\sigma\in\mathrm{Perm}(V)$ with $\lambda\ge0$); conversely every such combination is in $\mathrm{DS}(V,W)$ & \lean{GsimOT.exists_birkhoff_decomposition_pos}\newline{\scriptsize\texttt{Common/BirkhoffDecomposition.lean:72}} \newline \lean{GsimOT.exists_birkhoff_decomposition}\newline{\scriptsize\texttt{Common/BirkhoffDecomposition.lean:42}} \newline \lean{GsimOT.isDS_of_decomposition}\newline{\scriptsize\texttt{Common/BirkhoffDecomposition.lean:23}} \\ \midrule
{}Observation~21 & 20 & $L_Q(v,w)\ge|\deg_G(v)-\deg_H(w)|$ & \lean{GsimOT.observation21}\newline{\scriptsize\texttt{Common/Observation21\_DegreeBound.lean:21}} \newline \lean{GsimOT.observation21_bij}\newline{\scriptsize\texttt{Common/Observation21\_DegreeBound.lean:40}} \\ \midrule
{}Lemma~22 (construction) & 21 & Adding the clique $K$ ($m=2\Delta+2\tau+k(\tau+\delta)+1$) and the staircase & \lean{GsimOT.lockSize}\newline{\scriptsize\texttt{Common/Lemma22\_LockingMarkers.lean:36}} \newline \lean{GsimOT.stair}\newline{\scriptsize\texttt{Common/Lemma22\_LockingMarkers.lean:40}} \newline \lean{GsimOT.lockGraph}\newline{\scriptsize\texttt{Common/Lemma22\_LockingMarkers.lean:56}} \\ \midrule
{}Lemma~22 (2) & 21 & $Q(z_j,z_j)=1$ $\Rightarrow$ $C_1(Q')$ w.r.t.\ $(G',H')$ equals $C_1(Q)$ w.r.t.\ $(G,H)$ & \lean{GsimOT.extendId}\newline{\scriptsize\texttt{Common/Lemma22\_LockingMarkers.lean:260}} \newline \lean{GsimOT.lemma22_cost}\newline{\scriptsize\texttt{Common/Lemma22\_LockingMarkers.lean:351}} \\ \midrule
{}Lemma~22 (1), $\diamondsuit=\mathrm{ed}$ & 21 & $C_{\mathrm{ed}}(\pi)<\tau$ $\Rightarrow$ $\pi$ fixes all markers and $\pi(V)=V$; via $C_1(\pi)\le C_{\mathrm{ed}}(\pi)$ and the case $\diamondsuit=1$ & \lean{GsimOT.costOneBij_le_costEditBij}\newline{\scriptsize\texttt{Common/Lemma22\_LockingMarkersEdit.lean:26}} \newline \lean{GsimOT.lemma22_bij_edit}\newline{\scriptsize\texttt{Common/Lemma22\_LockingMarkersEdit.lean:41}} \\ \midrule
{}Lemma~22 (2), $\diamondsuit=\mathrm{ed}$ & 21 & $Q(z_j,z_j)=1$ $\Rightarrow$ $C_{\mathrm{ed}}(Q')$ w.r.t.\ $(G',H')$ equals $C_{\mathrm{ed}}(Q)$ w.r.t.\ $(G,H)$ & \lean{GsimOT.lemma22_cost_edit}\newline{\scriptsize\texttt{Common/Lemma22\_LockingMarkersEdit.lean:53}} \\
\end{longtable}

\section*{Notes on the formalisation}

\begin{enumerate}[leftmargin=*]
\item Proposition~5 states $d_{\mathrm{ed}}(G,H)\ge4N-2$ and $d^{\mathrm{OT}}_{\mathrm{ed}}(G,H)\le3N+1$ (\lean{proposition5_dEdit}); \lean{proposition5} is the form with the explicit witness $Q$. \lean{proposition5_delta_of_le} is the computation of the proof for every odd $N\ge3$ with $\frac{10}{3(3N+1)}\le\varepsilon$.
\item The common vertex set of $G'$ and $H'$ is \lean{Base N} $=$ \lean{Fin N} $\oplus$ \lean{Fin N}: \lean{Sum.inl c} is the core vertex $v_{2c}$ of $G'$ and $w_{2c+1}$ of $H'$ (the identification $v_i\leftrightarrow w_{i+1}$ of the paper), \lean{Sum.inr r} is the marker $z_r$, and $\mathbb Z_{2N}$ is $\{0,\dots,2N-1\}$. With this identification, $\pi_+$ is the identity on the core and $\pi_-$ maps the core vertex $c$ to $c-1\bmod N$ (\lean{predCore}).
\item A compliant $Q$ is \lean{extendId (extMarkers R)}, with $R$ the restriction of $Q$ to the core vertices. Lemma~23 is computed on $(G',H')$ (\lean{costEdit_extMarkers}) and transferred to $(G,H)$ by Lemma~22~(2) (\lean{lemma23}).
\item Lemma~24: \lean{costEdit_witness} computes the cost of the witness, $C_{\mathrm{ed}}(Q)=3N$, and \lean{lemma24} derives $d^{\mathrm{OT}}_{\mathrm{ed}}(G,H)\le3N+1$ from it.
\item Lemma~25: instead of replacing $\pi$ on the clique by the identity, the pairs involving the clique are dropped from $C_{\mathrm{ed}}(\pi)$ (\lean{costEditBij_restrict}). Claim~26 is stated with $2k$, the number of preserved ordered core edges (\lean{pres}), and $\sum_{v_i\in S}d_i=0$ is proved with the involution $c\mapsto N-c$ on $S$.
\item Lemma~22 for $\diamondsuit=\mathrm{ed}$ is in \texttt{Common/Lemma22\_LockingMarkersEdit.lean}; its first item follows from the case $\diamondsuit=1$, as $C_1(\pi)\le C_{\mathrm{ed}}(\pi)$. The parameters $k=N$, $\Delta=N+1$, $\tau=4N-1$, $\delta=1$ are those of the paper (\lean{mK}, \lean{GN}).
\end{enumerate}

\end{document}